% Generated draft: requires mathematical and source review.
\documentclass[11pt]{article}
\usepackage[margin=25mm]{geometry}
\usepackage{fontspec}
\setmainfont{Latin Modern Roman}
\usepackage{amsmath,amssymb,mathtools,unicode-math}
\setmathfont{Latin Modern Math}
\usepackage{enumitem,needspace,xurl,hyperref}
\setlength{\parindent}{0pt}
\setlength{\parskip}{5pt}
\setlength{\emergencystretch}{3em}
\newcommand{\N}{\mathbb{N}}
\newcommand{\Z}{\mathbb{Z}}
\newcommand{\Q}{\mathbb{Q}}
\newcommand{\R}{\mathbb{R}}
\newcommand{\C}{\mathbb{C}}
\newcommand{\F}{\mathbb{F}}
\newcommand{\A}{\mathbb{A}}
\newcommand{\PP}{\mathbb{P}}
\newcommand{\E}{\mathbb{E}}
\newcommand{\eps}{\varepsilon}
\newcommand{\dd}{\,\mathrm{d}}
\begin{document}
\section*{UnsolvedMath ID: 30002526}
\textbf{Source problem code:} catalogue entry
\subsection{Definitions and mathematical statement}
A group \(G\) is finitely presented if it admits a presentation \(G\cong\langle g_1,\ldots,g_m\mid R_1,\ldots,R_n\rangle\) with finitely many generators and relations.

For a precise contextual interpretation, let \(X\) be a reduced irreducible projective variety over \(\C\) of complex dimension two, and let \(X^{\mathrm{an}}\) denote its associated analytic space. Say that \(X\) has only normal-crossing singularities if, at every point \(p\), its analytic germ is isomorphic to the germ at the origin of
\[
\{(z_1,z_2,z_3)\in\C^3:z_1\cdots z_r=0\},\qquad r\in\{1,2,3\}.
\]
Here \(r=1\) describes a smooth point. This is a local analytic condition: local branches need not belong to distinct global irreducible components.

Under these conventions, the supplied question is: for every finitely presented group \(G\), does there exist such an \(X\), with only normal-crossing singularities, and a point \(x\in X^{\mathrm{an}}\), such that
\[
\pi_1(X^{\mathrm{an}},x)\cong G?
\]
The fundamental group here is the ordinary topological fundamental group. The question asks whether Whitney umbrellas can be eliminated from the realization described in the supplied background while retaining irreducibility, projectivity, dimension two, and realization of every finitely presented group.

The complex ground field, analytic topology, and local analytic meaning of normal crossings are contextual interpretations requiring confirmation: the supplied statement does not explicitly specify them. In particular, the local condition above must not silently be strengthened to a global simple-normal-crossing condition.
\subsection{Short English statement}
Can every finitely presented group be realized as the fundamental group of an irreducible projective surface whose only singularities are normal crossings, eliminating the Whitney umbrellas mentioned in the supplied background? The complex-analytic interpretation used here requires confirmation from the source.
\subsection{Sources}
\begin{itemize}
\item Supplied local UnsolvedMath registry record, permanent ID \(30002526\), source problem code \texttt{catalogue entry}, titled ``Eliminating Whitney Umbrellas in Projective Group Realization.'' The record belongs to the Oberwolfach Reports Open Problems collection and gives the proposed year as 2014. No individual proposer is identified. Its original, clean, and displayed statements agree; the supplied verification labels the statement exact and publication-ready, without independent verification here.
\item Supplied report citation: \emph{Mini-Workshop: Kähler Groups} (2014), \emph{Oberwolfach Reports}, pp.\ 531--553, DOI \texttt{10.4171/owr/2014/09}. \url{https://doi.org/10.4171/owr/2014/09}. This is also the supplied literature-source link labelled ``Oberwolfach Report 9, 2014.'' The background describes these reports as documenting mathematical workshops and their open-problem discussions.
\item Supplied literature-source label: ``Kapovich, 2013.'' \url{https://doi.org/10.1007/s00222-013-0453-4}. No article title or precise theorem reference is supplied.
\item Supplied literature assessment, checked on \textbf{2026-08-21}: the record labels the problem open and states that Kapovich's irreducible projective realization still uses Whitney umbrellas in the stated theorem, and that removing them is not established by the cited construction. This is supplied dated curation derived from the OWR material, not an independently verified literature assessment. The supplied background explicitly calls for independent verification of time-sensitive status and cited claims.
\end{itemize}
\end{document}
